\section{Motivation}
Token-based voting captures binary preferences on formal proposals but does little to surface intensity. After a decade of on-chain experimentation, DAOs still struggle to translate intent into timely, accountable action. Common pain points include:
\begin{itemize}[leftmargin=*]
  \item Endless discussion threads with unclear next steps.
  \item Committees that drift or concentrate power.
  \item Fragmented knowledge bases that deter newcomers.
  \item Social pressure that favors incumbents over fresh ideas.
\end{itemize}

By leveraging preference intensity through time-lock commitments, \signals{} constructs a dynamic idea board where the most strongly supported initiatives rise naturally, keeping attention on what matters and when. Concretely, each position requires locking real tokens for time; supporters own an ERC-721 claim on those tokens and must wait for expiry (or success) to redeem, making opportunity cost explicit.

To counteract last-minute pile-ons and apathetic equilibria, the protocol optionally distributes participation rewards. A shared \texttt{IncentivesPool} funds supporters of accepted initiatives based on when they committed support (earlier commitments earn higher weight via time buckets), while a \texttt{Bounties} module lets contributors escrow targeted rewards per initiative. Both mechanisms aim to reward early, accurate signaling and raise the cost of low-conviction behavior without distorting the core lock-and-decay semantics.
